% Generated draft: requires mathematical and source review.
\documentclass[11pt]{article}
\usepackage[margin=25mm]{geometry}
\usepackage{fontspec}
\setmainfont{Latin Modern Roman}
\usepackage{amsmath,amssymb,mathtools,unicode-math}
\setmathfont{Latin Modern Math}
\usepackage{enumitem,needspace,xurl,hyperref}
\setlength{\parindent}{0pt}
\setlength{\parskip}{5pt}
\setlength{\emergencystretch}{3em}
\newcommand{\N}{\mathbb{N}}
\newcommand{\Z}{\mathbb{Z}}
\newcommand{\Q}{\mathbb{Q}}
\newcommand{\R}{\mathbb{R}}
\newcommand{\C}{\mathbb{C}}
\newcommand{\F}{\mathbb{F}}
\newcommand{\A}{\mathbb{A}}
\newcommand{\PP}{\mathbb{P}}
\newcommand{\E}{\mathbb{E}}
\newcommand{\eps}{\varepsilon}
\newcommand{\dd}{\,\mathrm{d}}
\begin{document}
\section*{UnsolvedMath ID: 30005908}
\textbf{Source problem code:} catalogue entry
\subsection{Definitions and mathematical statement}
Let \(k\) be a field. In the supplied excerpt, a finite tensor \(k\)-category \((\mathcal C,\otimes,\mathbf 1)\) is a locally finite \(k\)-linear abelian category with finitely many isomorphism classes of simple objects, enough projectives, projective covers, an associative tensor product with simple unit \(\mathbf 1\), and left and right duals for every object. Here locally finite is understood in its usual sense: objects have finite length and morphism spaces are finite-dimensional over \(k\). No assumption that \(k\) is algebraically closed is supplied. Examples given are finite-dimensional representations of a finite group and finite-dimensional modules over a finite-dimensional Hopf algebra.

Define the graded cohomology algebra
\[
 A=H^*(\mathcal C):=\bigoplus_{n\geq 0}\operatorname{Ext}_{\mathcal C}^n(\mathbf 1,\mathbf 1),
\]
with Yoneda multiplication. For an object \(X\in\mathcal C\), put
\[
 E(X):=\bigoplus_{n\geq 0}\operatorname{Ext}_{\mathcal C}^n(\mathbf 1,X),
\]
with its natural graded \(A\)-module structure induced by Yoneda composition.

\textbf{Finite-generation question.} A conventional precise formulation of the finite-generation question in the record is: for every field \(k\) and every finite tensor \(k\)-category \(\mathcal C\) as above, is \(A\) finitely generated as a \(k\)-algebra, and is \(E(X)\) finitely generated as a graded \(A\)-module for every \(X\in\mathcal C\)? The supplied curation identifies the general finite-generation problem as the Etingof--Ostrik conjecture. However, the excerpt says only that the cohomology of \(\mathcal C\) is finitely generated over \(A\); it does not specify the coefficient Ext-groups or their quantifiers. Thus the displayed coefficient condition is an explicitly stated conventional interpretation requiring confirmation against the cited source.

\textbf{Perfect complexes and the second question.} Conventionally, a perfect complex over \(\mathcal C\) is an object of the bounded derived category \(D^b(\mathcal C)\) represented by a bounded complex of projective objects. For an object \(X\), write \(\ell(X)\) for its composition length and \(\operatorname{FPdim}_{\mathcal C}(X)\) for its Frobenius--Perron dimension. For a perfect complex \(D\), define
\[
 L(D)=\sum_{n\in\mathbb Z}\ell(H^n(D)),\qquad
 F(D)=\sum_{n\in\mathbb Z}\operatorname{FPdim}_{\mathcal C}(H^n(D)).
\]
These sums are finite. The record also asks what homological lower bounds analogous to Carlsson's conjecture hold for perfect complexes. It supplies no definite proposed inequality or complete hypotheses for this question, so this part remains a research question rather than a uniquely specified conjecture.

\textbf{Supplied contextual result, with unresolved transcription.} The excerpt attributes to its reference [1] a theorem for a field of characteristic different from \(2\) and a finite tensor category with finitely generated cohomology: if \(d=\dim_{\mathrm{Krull}}H^*(\mathcal C)\geq 8\), there are infinitely many pairwise non-isomorphic, nontrivial perfect complexes \(D\) satisfying an inequality apparently intended to read
\[
 L(D)\leq F(D)<2^{d}.
\]
The extracted formula has displaced summation symbols and the terminal text \(2d\), without an explicit exponent. The displayed reconstruction, including the exponent, therefore requires source confirmation. The excerpt also does not specify whether non-isomorphism is in a category of complexes, a homotopy category, or a derived category; under the conventional derived interpretation, nontrivial means that some \(H^n(D)\) is nonzero. Its convention for the Krull dimension of the graded cohomology algebra is not stated.

The excerpt separately attributes to Jon Carlson, in reference [2], an analogous result for every finite group \(G\) over a field \(k\) of odd characteristic \(p\), when its \(p\)-rank \(d\) is at least \(8\). Here the \(p\)-rank is the largest integer \(r\) such that \(G\) contains a subgroup isomorphic to \((\mathbb Z/p\mathbb Z)^r\). The claimed conclusion is infinitely many pairwise non-isomorphic, nontrivial perfect complexes with the provisionally reconstructed bound
\[
 \sum_{n\in\mathbb Z}\dim_k H^n(D)<2^{d}.
\]
The same missing-exponent qualification applies. These supplied results are background claims, not assertions of new results or independently verified theorems.

The historical discussion distinguishes algebraic and topological versions of Carlsson's conjecture, attributes an implication from the algebraic version to the topological version to reference [3], and says that Carlsson settled the case \(p=2\), \(d\leq 3\), in reference [4]. It attributes odd-characteristic counterexamples for \(d\geq 8\) to Srikanth Iyengar and Mark Walker in reference [6], using exterior algebras and Lefschetz elements. The relevant conjectures, hypotheses, and bibliographic entries are absent, so these historical statements cannot supply a precise additional conjecture here.
\subsection{Short English statement}
The main question asks whether every finite tensor category has a finitely generated cohomology algebra and finitely generated cohomology with coefficients in every object. The record also asks about lower bounds for the homology of perfect complexes, but leaves that target unspecified and contains damaged formulas in its contextual results. This is a statement-only draft requiring mathematical and source review.
\subsection{Sources}
\begin{itemize}
\item Local registry: \texttt{data/unsolvedmath/problems.json}, with \texttt{catalog/problems.json} as the supplementary registry; permanent UnsolvedMath ID \(30005908\), exact source problem code \texttt{catalogue entry}. The supplied record is titled \emph{Cohomological Finite Generation for Finite Tensor Categories}, with proposed year 2024 and no named proposer.
\item Supplied report citation: \emph{Hochschild (Co)Homology and Applications} (2024), Oberwolfach Reports, pp.\ 1097--1162, DOI 10.4171/owr/2024/20. \url{https://doi.org/10.4171/owr/2024/20}. This is the report-level citation supplied for the excerpt; the particular contribution, author, and exact problem page are not identified.
\item Supplied literature source: Nakano, Vashaw and Yakimov, \url{https://arxiv.org/abs/2209.11621}. The record supplies no title or explicit identification of this source with any numbered reference in the excerpt.
\item Supplied literature source: Oberwolfach Report 20/2024, \url{https://doi.org/10.4171/owr/2024/20}.
\item Literature assessment dated \texttt{2026-08-22}: the supplied curation labels the problem partially solved and states that finite generation is known for several classes of finite tensor categories while the arbitrary-category Etingof--Ostrik conjecture remains unresolved. This is supplied dated curation, not an independently verified literature assessment. The background expressly requests independent verification of time-sensitive status and cited claims.
\item The excerpt cites references [1]--[6] but supplies no corresponding bibliography. Its attributions include Etingof--Ostrik through the curation, Jon Carlson, Carlsson, and Srikanth Iyengar and Mark Walker; no additional bibliographic identities are inferred.
\end{itemize}
\end{document}
